% ID: SEL-D0083-OA-E01
% ESTADO: cotejada. Segunda pasada visual y borrador compilado revisados.
% Fuente DOC: Selectividad/sel_2012_matematicas_aplicadas/5 MATEMÁT. APLIC. CC. SS/EXAMEN 3-ANDALUCÍA.doc
% SHA-256 original: c8688a19515dfeb348e8f8817d3f4a6c0a967723c270f2ec42e7f365d001e555
% Página convertida: 1; opción A; ejercicio 1.

\textbf{(2.5 puntos)} Un empresario fabrica camisas y pantalones para jóvenes. Para hacer una camisa se necesitan 2 metros de tela y 5 botones, y para hacer un pantalón hacen falta 3 metros de tela, 2 botones y 1 cremallera. La empresa dispone de 1050 metros de tela, 1250 botones y 300 cremalleras. El beneficio que se obtiene por la venta de una camisa es de 30 euros y el de un pantalón es de 50 euros.

Suponiendo que se vende todo lo que se fabrica, calcule el número de camisas y de pantalones que debe confeccionar para obtener el máximo beneficio, y determine este beneficio máximo.
